\chapter{Configurations et preuves géométriques}
\section*{Exercice 11.1 -- Milieux}
Dans un triangle $ABC$, $M$ et $N$ sont les milieux de $[AB]$ et $[AC]$. Démontrer $\overrightarrow{MN}=\tfrac12\overrightarrow{BC}$.

\section*{Exercice 11.2 -- Barycentre}
Approfondissement : trouver le barycentre de $A(0,2)$ de poids $1$ et $B(6,-1)$ de poids $2$.

\section*{Exercice 11.3 -- Médianes}
Approfondissement : pour $A(0,0)$, $B(6,0)$, $C(0,9)$, trouver l'intersection des médianes et sa position sur la médiane issue de $A$.

\section*{Exercice 11.4 -- Cercle circonscrit}
Approfondissement : pour le triangle de l'exercice 7.2, trouver le centre et le rayon du cercle circonscrit.

\section*{Exercice 11.5 -- Aire}
Approfondissement : deux côtés d'un triangle mesurent $5$ et $8$, et leur angle compris est $30\degre$. Calculer l'aire et expliquer la formule utilisée.

\section*{Exercice 11.6 -- Lieu}
Approfondissement : déterminer le lieu des points équidistants de $F(0,2)$ et de l'axe des abscisses. Vérifier le point $(2,2)$.
